\section*{Chapitre \ref{ch:hf:latency-arbitrage-and-its-defence} --- Arbitrage de latence et sa défense}

\begin{solution}{exo:hf:latency-arbitrage-and-its-defence:1}
$2-0.5=1.5$ ticks par action, avant frais.
\end{solution}

\begin{solution}{exo:hf:latency-arbitrage-and-its-defence:2}
$537/510\approx1.05$ par minute.
\end{solution}

\begin{solution}{exo:hf:latency-arbitrage-and-its-defence:3}
Plus de $60-45=15$ microsecondes.
\end{solution}

\begin{solution}{exo:hf:latency-arbitrage-and-its-defence:4}
Voir la \cref{prop:hf:latency-arbitrage-and-its-defence:bump}~: les instants d'arrivée sont décalés de la même $d$ et gardent leur ordre. Un délai symétrique, appliqué aussi
aux données sortantes de la plateforme, empêche encore les traders de réagir aux exécutions de la plateforme plus vite que celle-ci ne peut acheminer l'ordre d'un client.
\end{solution}

\begin{solution}{exo:hf:latency-arbitrage-and-its-defence:5}
Lorsque l'événement survient juste avant la fin d'un lot, un preneur de liquidité peut arriver avant la frontière tandis que l'annulation, plus lente, arrive après elle~: le preneur se trouve dans
un lot antérieur. La probabilité vaut environ l'écart entre les deux latences divisé par l'intervalle, soit ici environ 2~\%.
\end{solution}

\begin{solution}{exo:hf:latency-arbitrage-and-its-defence:6}
À mi-chemin entre 28.8~\% et 13.4~\%, on trouve 21.1~\%~: le garde-fou y parvient encore à 100~ms (18.7~\%) et à 300~ms (20.0~\%), mais pas à 600~ms (24.6~\%).
\end{solution}

\begin{solution}{exo:hf:latency-arbitrage-and-its-defence:7}
0.314. La probabilité tombe à 1.0~\% pour un délai de 110~microsecondes et à 0.5~\% pour 120~: le délai doit compenser le désavantage du fournisseur de liquidité et la majeure partie de
la gigue.
\end{solution}

\begin{solution}{exo:hf:latency-arbitrage-and-its-defence:8}
Un flux de carnet d'ordres ne montre que les ordres qui ont abouti~; les perdants d'une course (ordres manqués, annulations arrivées trop tard) ne laissent aucune trace dans ce flux.
Les courses ne sont visibles que dans les données de messages qui contiennent les tentatives échouées, comme l'ont fait Aquilina, Budish et O'Neill, ou dans ses propres journaux d'annulations rejetées et
des ordres preneurs manqués.
\end{solution}

\begin{solution}{pb:hf:latency-arbitrage-and-its-defence:1}
\begin{enumerate}
\item Voir la \cref{def:hf:latency-arbitrage-and-its-defence:arb}~; l'information est publique et seule la vitesse compte, contrairement au trading informé, qui exploite
une information que les autres n'ont pas.
\item La durée médiane de l'opportunité est passée de 97 à 7 millisecondes~; son profit médian est resté proche de 0.08 point d'indice.
\item Environ 537 courses par jour et par valeur du FTSE 100~; écart modal de 5--10 microsecondes~; environ 22~\% du volume en courses~; taxe de 0.42 points de base sur l'ensemble du volume.
\item Les données de messages contiennent les ordres et annulations échoués, c'est-à-dire les perdants~; un flux de carnet d'ordres ne montre que ce qui a abouti.
\item Une annulation du fournisseur de liquidité, à une médiane de 50 microsecondes, contre trois snipers à 40, gigue lognormale 0.3~; un saut d'un tick contre une cotation à demi-tick.
\item 0.899 et 0.449 ticks par action et par course.
\item 12~\%, 75~\% et 99.9~\%.
\item Une médiane de 7.3 microsecondes, conforme aux 5--10 microsecondes rapportées pour Londres.
\item Voir la \cref{def:hf:latency-arbitrage-and-its-defence:bump}.
\item Voir la preuve~: un décalage commun conserve l'ordre~; si seuls les preneurs sont retardés, la comparaison devient $x_{\text{prise}}+d<x_{\text{annulation}}$.
\item 0.899 en continu et avec un délai symétrique de 350~microsecondes~; 0 avec le même délai appliqué aux seuls preneurs~; 0.211 et 0.021 pour des lots de 100
microsecondes et d'une milliseconde.
\item L'IEX a été approuvé avec un délai de 350 microsecondes (2016) et son ordre D-Limit repricé par la plateforme (2020)~; le délai asymétrique de quatre millisecondes de l'EDGA
a été désapprouvé (2020)~; l'abrogation de la règle de trade-through a été proposée (2026).
\item Le fournisseur de liquidité peut rejeter une demande de transaction après un temps de rétention si le cours a bougé~; il le contrôle, dans le cadre des
principes du FX Global Code.
\item 90~\% de la demi-fourchette dans un carnet continu et avec un délai symétrique ($0.449/0.5$), zéro avec un délai asymétrique, 2.1~\% avec un
lot d'une milliseconde.
\item 13.4~\% et 0.402 à 1~ms, 20.0~\% et 0.330 à 300~ms, 24.6~\% et 0.192 à 600~ms~: le garde-fou fonctionne tant qu'il lit le meneur plus vite que ne
dure l'avance, et échoue au-delà.
\item 31~\% du volume à 1~ms (5\,225 actions par séance contre 7\,550), et de messages.
\item Un garde-fou alimenté aussi vite que ses moyens le permettent, une cotation légèrement plus large, et une préférence pour les plateformes qui protègent les cotations.
\item Le délai rend leurs exécutions plus lentes et moins certaines~; les routeurs qui valorisent la vitesse d'exécution vont ailleurs, comme
le montre le registre des objections de la SEC.
\item Les plateformes ne seraient plus tenues d'acheminer les ordres vers les cotations protégées des autres, ni de les respecter~; les cotations périmées pourraient faire l'objet de trade-through au lieu d'être prises,
modifiant ainsi les courses qui valent la peine d'être disputées.
\item Les fournisseurs de liquidité et, par leurs fourchettes plus larges, tous ceux qui négocient.
\end{enumerate}
\end{solution}

\begin{solution}{iq:hf:latency-arbitrage-and-its-defence:1}
Sélection adverse~: la contrepartie savait quelque chose que l'on ignorait. Sniping~: l'information était publique, et la contrepartie a été plus rapide
à agir dessus que l'on ne l'était pour annuler.
\iqlookfor{information privée contre information publique.}
\end{solution}

\begin{solution}{iq:hf:latency-arbitrage-and-its-defence:2}
Consigner chaque annulation arrivée trop tard (rejets, exécutions entre l'envoi et l'accusé de réception) et les exécutions qui ont suivi un mouvement sur une autre plateforme dans
l'aller-retour de l'ordre~; comparer les horodatages du mouvement, de l'annulation et de l'exécution.
\iqlookfor{annulations tardives et horodatages entre plateformes.}
\end{solution}

\begin{solution}{iq:hf:latency-arbitrage-and-its-defence:3}
Le chemin depuis le paquet de données de marché jusqu'à l'annulation sur le fil~: kernel bypass, décodage du flux, décision, encodage de l'ordre, carte réseau, la
liaison vers la bourse~; mesurer chaque étape avant d'optimiser (livre~13, chapitre~1).
\iqlookfor{une vue mesurée, étape par étape.}
\end{solution}

\begin{solution}{iq:hf:latency-arbitrage-and-its-defence:4}
Dans un lot, tous les ordres sont traités de la même façon et ont le même prix, de sorte qu'arriver plus tôt ne rapporte rien~; à une frontière, arriver juste avant plutôt que
juste après, on place l'ordre dans une enchère antérieure.
\iqlookfor{temps discret et ses frontières.}
\end{solution}

\begin{solution}{iq:hf:latency-arbitrage-and-its-defence:5}
La maison devient plus lente que ses concurrents sur cette route~: ses cotations y sont périmées plus longtemps et ses stratégies preneuses perdent des courses. Élargir ou
retirer les cotations concernées, arrêter les stratégies sensibles à la latence, et revenir à la route par fibre, avec sa latence dans les paramètres de risque.
\iqlookfor{dépendance opérationnelle à la vitesse et solution de repli.}
\end{solution}

\begin{solution}{iq:hf:latency-arbitrage-and-its-defence:6}
Le minimum de $n$ exponentielles de moyenne $s$ est une exponentielle de paramètre $n/s$~; il bat une exponentielle indépendante de paramètre $1/c$ avec probabilité
$(n/s)/(n/s+1/c)=nc/(nc+s)$.
\iqlookfor{exponentielles en compétition.}
\end{solution}
